\documentclass[12pt, xcolor=dvipsnames,svgnames,x11names]{article}

\usepackage{tikz}
\usepackage{pgfplots}
\usepackage{minted}





\def\m@th{\normalcolor\mathsurround\z@}

\setlength{\parskip}{\baselineskip}%
\setlength{\parindent}{0pt}%

\def\e{{\textrm{e}}}
\def\d{\textrm{d}}
\def\T{{\textsf{T}}}
\def\sech{{\textrm{sech}}}
\def\x{{\mathbf x}}

\def\++{{$^{++}$}}

\def\Abb{{\mathbb A}}
\def\Bbb{{\mathbb B}}
\def\Cbb{{\mathbb C}}
\def\Dbb{{\mathbb D}}
\def\Ebb{{\mathbb E}}
\def\Fbb{{\mathbb F}}
\def\Gbb{{\mathbb G}}
\def\Hbb{{\mathbb H}}
\def\Ibb{{\mathbb I}}
\def\Jbb{{\mathbb J}}
\def\Kbb{{\mathbb K}}
\def\Lbb{{\mathbb L}}
\def\Mbb{{\mathbb M}}
\def\Nbb{{\mathbb N}}
\def\Obb{{\mathbb O}}
\def\Pbb{{\mathbb P}}
\def\Qbb{{\mathbb Q}}
\def\Rbb{{\mathbb R}}
\def\Sbb{{\mathbb S}}
\def\Tbb{{\mathbb T}}
\def\Ubb{{\mathbb U}}
\def\Vbb{{\mathbb V}}
\def\Wbb{{\mathbb W}}
\def\Xbb{{\mathbb X}}
\def\Ybb{{\mathbb Y}}
\def\Zbb{{\mathbb Z}}



\def\Acal{{\mathcal A}}
\def\Bcal{{\mathcal B}}
\def\Ccal{{\mathcal C}}
\def\Dcal{{\mathcal D}}
\def\Ecal{{\mathcal E}}
\def\Fcal{{\mathcal F}}
\def\Gcal{{\mathcal G}}
\def\Hcal{{\mathcal H}}
\def\Ical{{\mathcal I}}
\def\Jcal{{\mathcal J}}
\def\Kcal{{\mathcal K}}
\def\Lcal{{\mathcal L}}
\def\Mcal{{\mathcal M}}
\def\Ncal{{\mathcal N}}
\def\Ocal{{\mathcal O}}
\def\Pcal{{\mathcal P}}
\def\Qcal{{\mathcal Q}}
\def\Rcal{{\mathcal R}}
\def\Scal{{\mathcal S}}
\def\Tcal{{\mathcal T}}
\def\Ucal{{\mathcal U}}
\def\Vcal{{\mathcal V}}
\def\Wcal{{\mathcal W}}
\def\Xcal{{\mathcal X}}
\def\Ycal{{\mathcal Y}}
\def\Zcal{{\mathcal Z}}

\def\abf{{\mathbf a}}
\def\bbf{{\mathbf b}}
\def\cbf{{\mathbf c}}
\def\dbf{{\mathbf d}}
\def\ebf{{\mathbf e}}
\def\fbf{{\mathbf f}}
\def\gbf{{\mathbf g}}
\def\hbf{{\mathbf h}}
\def\ibf{{\mathbf i}}
\def\jbf{{\mathbf j}}
\def\kbf{{\mathbf k}}
\def\lbf{{\mathbf l}}
\def\mbf{{\mathbf m}}
\def\nbf{{\mathbf n}}
\def\obf{{\mathbf o}}
\def\pbf{{\mathbf p}}
\def\qbf{{\mathbf q}}
\def\rbf{{\mathbf r}}
\def\sbf{{\mathbf s}}
\def\tbf{{\mathbf t}}
\def\ubf{{\mathbf u}}
\def\vbf{{\mathbf v}}
\def\wbf{{\mathbf w}}
\def\xbf{{\mathbf x}}
\def\ybf{{\mathbf y}}
\def\zbf{{\mathbf z}}

\def\Abf{{\mathbf A}}
\def\Bbf{{\mathbf B}}
\def\Cbf{{\mathbf C}}
\def\Dbf{{\mathbf D}}
\def\Ebf{{\mathbf E}}
\def\Fbf{{\mathbf F}}
\def\Gbf{{\mathbf G}}
\def\Hbf{{\mathbf H}}
\def\Ibf{{\mathbf I}}
\def\Jbf{{\mathbf J}}
\def\Kbf{{\mathbf K}}
\def\Lbf{{\mathbf L}}
\def\Mbf{{\mathbf M}}
\def\Nbf{{\mathbf N}}
\def\Obf{{\mathbf O}}
\def\Pbf{{\mathbf P}}
\def\Qbf{{\mathbf Q}}
\def\Rbf{{\mathbf R}}
\def\Sbf{{\mathbf S}}
\def\Tbf{{\mathbf T}}
\def\Ubf{{\mathbf U}}
\def\Vbf{{\mathbf V}}
\def\Wbf{{\mathbf W}}
\def\Xbf{{\mathbf X}}
\def\Ybf{{\mathbf Y}}
\def\Zbf{{\mathbf Z}}

\def\phat{{\hat p}}

\usepackage{amsmath}
\DeclareMathOperator*{\argmax}{arg\,max}
\DeclareMathOperator*{\argmin}{arg\,min}

\def\xbar{{\overline x}}
\def\ybar{{\overline y}}
\def\e{{\textrm{e}}}
\def\d{\textrm{d}}
\def\T{{\textsf{T}}}
\def\sech{{\textrm{sech}}}
\def\x{{\mathbf x}}

\def\++{{$^{++}$}}

\def\Abb{{\mathbb A}}
\def\Bbb{{\mathbb B}}
\def\Cbb{{\mathbb C}}
\def\Dbb{{\mathbb D}}
\def\Ebb{{\mathbb E}}
\def\Fbb{{\mathbb F}}
\def\Gbb{{\mathbb G}}
\def\Hbb{{\mathbb H}}
\def\Ibb{{\mathbb I}}
\def\Jbb{{\mathbb J}}
\def\Kbb{{\mathbb K}}
\def\Lbb{{\mathbb L}}
\def\Mbb{{\mathbb M}}
\def\Nbb{{\mathbb N}}
\def\Obb{{\mathbb O}}
\def\Pbb{{\mathbb P}}
\def\Qbb{{\mathbb Q}}
\def\Rbb{{\mathbb R}}
\def\Sbb{{\mathbb S}}
\def\Tbb{{\mathbb T}}
\def\Ubb{{\mathbb U}}
\def\Vbb{{\mathbb V}}
\def\Wbb{{\mathbb W}}
\def\Xbb{{\mathbb X}}
\def\Ybb{{\mathbb Y}}
\def\Zbb{{\mathbb Z}}



\def\Acal{{\mathcal A}}
\def\Bcal{{\mathcal B}}
\def\Ccal{{\mathcal C}}
\def\Dcal{{\mathcal D}}
\def\Ecal{{\mathcal E}}
\def\Fcal{{\mathcal F}}
\def\Gcal{{\mathcal G}}
\def\Hcal{{\mathcal H}}
\def\Ical{{\mathcal I}}
\def\Jcal{{\mathcal J}}
\def\Kcal{{\mathcal K}}
\def\Lcal{{\mathcal L}}
\def\Mcal{{\mathcal M}}
\def\Ncal{{\mathcal N}}
\def\Ocal{{\mathcal O}}
\def\Pcal{{\mathcal P}}
\def\Qcal{{\mathcal Q}}
\def\Rcal{{\mathcal R}}
\def\Scal{{\mathcal S}}
\def\Tcal{{\mathcal T}}
\def\Ucal{{\mathcal U}}
\def\Vcal{{\mathcal V}}
\def\Wcal{{\mathcal W}}
\def\Xcal{{\mathcal X}}
\def\Ycal{{\mathcal Y}}
\def\Zcal{{\mathcal Z}}

\def\abf{{\mathbf a}}
\def\bbf{{\mathbf b}}
\def\cbf{{\mathbf c}}
\def\dbf{{\mathbf d}}
\def\ebf{{\mathbf e}}
\def\fbf{{\mathbf f}}
\def\gbf{{\mathbf g}}
\def\hbf{{\mathbf h}}
\def\ibf{{\mathbf i}}
\def\jbf{{\mathbf j}}
\def\kbf{{\mathbf k}}
\def\lbf{{\mathbf l}}
\def\mbf{{\mathbf m}}
\def\nbf{{\mathbf n}}
\def\obf{{\mathbf o}}
\def\pbf{{\mathbf p}}
\def\qbf{{\mathbf q}}
\def\rbf{{\mathbf r}}
\def\sbf{{\mathbf s}}
\def\tbf{{\mathbf t}}
\def\ubf{{\mathbf u}}
\def\vbf{{\mathbf v}}
\def\wbf{{\mathbf w}}
\def\xbf{{\mathbf x}}
\def\ybf{{\mathbf y}}
\def\zbf{{\mathbf z}}

\def\Abf{{\mathbf A}}
\def\Bbf{{\mathbf B}}
\def\Cbf{{\mathbf C}}
\def\Dbf{{\mathbf D}}
\def\Ebf{{\mathbf E}}
\def\Fbf{{\mathbf F}}
\def\Gbf{{\mathbf G}}
\def\Hbf{{\mathbf H}}
\def\Ibf{{\mathbf I}}
\def\Jbf{{\mathbf J}}
\def\Kbf{{\mathbf K}}
\def\Lbf{{\mathbf L}}
\def\Mbf{{\mathbf M}}
\def\Nbf{{\mathbf N}}
\def\Obf{{\mathbf O}}
\def\Pbf{{\mathbf P}}
\def\Qbf{{\mathbf Q}}
\def\Rbf{{\mathbf R}}
\def\Sbf{{\mathbf S}}
\def\Tbf{{\mathbf T}}
\def\Ubf{{\mathbf U}}
\def\Vbf{{\mathbf V}}
\def\Wbf{{\mathbf W}}
\def\Xbf{{\mathbf X}}
\def\Ybf{{\mathbf Y}}
\def\Zbf{{\mathbf Z}}

\def\phat{{\hat p}}
\def\what{{\hat w}}
\def\yhat{{\hat y}}


\newcommand{\answer}{
    \fcolorbox{orange}{PapayaWhip}{
\begin{minipage}{\textwidth}
        \textbf{Answer}
    \end{minipage}
}
}

\newcommand{\codeoutput}[1]{
Output:\\\vspace{5px}
    \fcolorbox{orange}{WhiteSmoke}{
\begin{minipage}{\textwidth}
        { \color{DodgerBlue} 
        
        
       \texttt{#1}
       }
    \end{minipage}
}
}


\newcommand{\orangebox}[1]{
    \fcolorbox{orange}{white}{
\begin{minipage}{\textwidth}
       {\color{MidnightBlue} 
       #1
       }
    \end{minipage}
}
}


\definecolor{lightapricot}{rgb}{0.99, 0.84, 0.69}
\definecolor{lightblue}{rgb}{0.68, 0.85, 0.90}
\definecolor{blanchedalmond}{rgb}{1.0, 0.92, 0.8}
\definecolor{blizzardblue}{rgb}{0.67, 0.9, 0.93}
\definecolor{blond}{rgb}{0.98, 0.94, 0.75}
\definecolor{babyblueeyes}{rgb}{0.63, 0.79, 0.95}
\definecolor{babypink}{rgb}{0.96, 0.76, 0.76}
\definecolor{bananamania}{rgb}{0.98, 0.91, 0.71}
\definecolor{beaublue}{rgb}{0.74, 0.83, 0.9}
\definecolor{antiquewhite}{rgb}{0.98, 0.92, 0.84}
\definecolor{anti-flashwhite}{rgb}{0.95, 0.95, 0.96}
\definecolor{brightlavender}{rgb}{0.75, 0.58, 0.89}
\definecolor{brightube}{rgb}{0.82, 0.62, 0.91}
\definecolor{brilliantlavender}{rgb}{0.96, 0.73, 1.0}
\definecolor{lightcerulean}{rgb}{0.11, 0.67, 0.84}
\definecolor{cpeyellow}{HTML}{F5F5DC}                 % Beige (much softer)
\definecolor{powderGreen}{HTML}{F0F8E8}               % Very pale mint


% Darker versions of the original colors
\definecolor{deepapricot}{rgb}{0.85, 0.65, 0.45}        % darker lightapricot
\definecolor{deepblue}{rgb}{0.45, 0.65, 0.75}           % darker lightblue
\definecolor{toastedalmond}{rgb}{0.85, 0.75, 0.60}      % darker blanchedalmond
\definecolor{stormblue}{rgb}{0.45, 0.70, 0.75}          % darker blizzardblue
\definecolor{goldenrod}{rgb}{0.80, 0.75, 0.50}          % darker blond
\definecolor{royalblue}{rgb}{0.40, 0.60, 0.80}          % darker babyblueeyes
\definecolor{rosewood}{rgb}{0.75, 0.45, 0.45}           % darker babypink
\definecolor{mustardyellow}{rgb}{0.80, 0.70, 0.45}      % darker bananamania
\definecolor{navyblue}{rgb}{0.50, 0.60, 0.70}           % darker beaublue
\definecolor{antiquetan}{rgb}{0.80, 0.70, 0.60}         % darker antiquewhite
\definecolor{silvergray}{rgb}{0.75, 0.75, 0.80}         % darker anti-flashwhite
\definecolor{darklavender}{rgb}{0.55, 0.35, 0.70}       % darker brightlavender
\definecolor{deepube}{rgb}{0.65, 0.40, 0.75}            % darker brightube
\definecolor{plumviolet}{rgb}{0.75, 0.50, 0.85}         % darker brilliantlavender
\definecolor{deepcerulean}{rgb}{0.08, 0.45, 0.65}       % darker lightcerulean

\definecolor{roseRed}{HTML}{e20047}
\definecolor{coolGray}{HTML}{474747} % Neutral gray to contrast roseRed
\definecolor{mocha}{HTML}{a47864}
\definecolor{sage}{HTML}{8a9a5b}
\definecolor{dustyblue}{HTML}{6e8ca0}
\definecolor{terracotta}{HTML}{c87f5b}
\definecolor{lavender}{HTML}{9d8bb0}
\definecolor{darkmocha}{HTML}{7a5a4b}
\definecolor{darksage}{HTML}{667244}
\definecolor{darkdustyblue}{HTML}{526878}
\definecolor{darkterracotta}{HTML}{9e6347}
\definecolor{darklavender}{HTML}{766688}
\definecolor{winery}{HTML}{7e212a}

% ---- Title header palette ----
\definecolor{hdrplum}{HTML}{2E0014}    % near-black plum, header band
\definecolor{hdrcoral}{HTML}{E17564}   % coral, accent rules
\definecolor{hdrslate}{HTML}{2F4F4F}   % dark slate grey, frame and body text
\colorlet{hdrslatetint}{hdrslate!12!white}
\colorlet{hdrcoraltint}{hdrcoral!15!white}

\usepackage{sectsty}
\sectionfont{\color{darklavender}}  % sets colour of sections
\subsectionfont{\color{plumviolet}}  % sets colour of sections
\subsubsectionfont{\color{deepube}}  % sets colour of sections

\usepackage{xcolor}
%\usepackage{universityos}
\usepackage{xspace}
\usepackage{url}
\usepackage[colorlinks=true,bookmarks=false,linkcolor=black,urlcolor=blue,citecolor=black]{hyperref}
\usepackage{fancyhdr}
\usepackage[top=1in,bottom=1in,left=1in,right=1in]{geometry}
\usepackage{multicol}
\usepackage{pgfplots}
\usetikzlibrary{circuits.logic.US}
% ---- Typography: Palatino for text and math ----
% The previous version loaded mathpazo and then restored the old \rmdefault,
% so Palatino applied to math only. The lines below apply it to the text too.
\usepackage[T1]{fontenc}
\usepackage[sc]{mathpazo}   % Palatino text + math, with real small caps
\linespread{1.05}           % Palatino has a large x-height, so open up the leading
\usepackage[scaled=0.92]{helvet}  % companion sans, used by \textsf
\renewcommand{\ttdefault}{cmtt}    % keep Computer Modern Typewriter for code

\usepackage{eso-pic}



\usetikzlibrary{calc}
\usepackage{circuitikz}



\usepackage[many]{tcolorbox}
\setlength{\parindent}{0pt}
\setlength{\parskip}{4pt}%

\newcommand{\coursenumber}{CPE 486/586}
\newcommand{\courseyear}{2026\xspace}
\newcommand{\coursedatetime}{Tue/Thur 11:20 AM -- 12:40 PM\xspace}
\newcommand{\coursetitle}{Machine Learning for Engineering Applications\xspace}
\newcommand{\courseterm}{Fall, \courseyear}
\newcommand{\instructorname}{Rahul Bhadani\xspace}
\newcommand{\instructoremail}{\href{mailto:rahul.bhadani@uah.edu}{rahul.bhadani@uah.edu}}
\newcommand{\instructoroffice}{ENG 217-H\xspace}

\renewcommand{\familydefault}{\rmdefault}   % Palatino (loaded in common.tex)


% Redefine \maketitle to include tcolorbox
% Palette: hdrplum (2E0014), hdrcoral (E17564), hdrslate (2F4F4F)
\makeatletter
\newcommand{\hdrrule}{{\color{hdrcoral}\rule{4.5em}{1.1pt}}}
\renewcommand{\maketitle}{%
      \begin{center}
         \begin{tcolorbox}[
               enhanced,
               colback=white,
               colframe=hdrslate,
               boxrule=0.8pt,
               arc=4pt,
               left=0pt,
               right=0pt,
               top=0pt,
               bottom=0pt,
               drop shadow={hdrslate!35!white},
         ]
               % Header band: deep plum with a coral baseline
               \begin{tcolorbox}[
                  enhanced,
                  colback=hdrplum,
                  colframe=hdrplum,
                  boxrule=0pt,
                  arc=4pt,
                  sharp corners=south,
                  left=2em,
                  right=2em,
                  top=1.1em,
                  bottom=1.1em,
                  borderline south={2.2pt}{0pt}{hdrcoral},
               ]
                  \centering
                  % Palatino has no bold small caps, so the title stays at medium
                  % weight and gains presence from size and the dark band instead.
                  {\LARGE\color{white}\linespread{1.25}\selectfont\@title}
               \end{tcolorbox}

               % Body: author, coral divider, then the due information
               \vspace{1.1em}
               \centering
               {\large\color{hdrslate}\@author}\\[0.7em]
               \hdrrule\\[0.7em]
               {\normalsize\color{hdrplum}\@date}
               \vspace{1.1em}
         \end{tcolorbox}
      \end{center}
      \vspace{1.5em}
}
\makeatother




\newcommand{\hw}{Homework 3:\\\normalsize Optimization and Linear Regression}
\newcommand{\duedate}{October 11, 2026, 11:59 PM}

\newenvironment{multiequation}[1][]{%
\begin{equation}%
   \begin{aligned}%
   \ifx#1\@empty\else\label{#1}\fi%
}{%
   \end{aligned}%
\end{equation}%
}
\pagestyle{fancy}

\lhead{\footnotesize{\coursenumber, \courseterm}, \footnotesize{UAH}}
\chead{}
\rhead{\footnotesize{\emph{\hw}}}
\lfoot{\footnotesize{\instructorname}}
\cfoot{\footnotesize{\thepage}}
\rfoot{\footnotesize{\textit{Last Revised:~\today}}}
\renewcommand{\headrulewidth}{0.1pt}
\renewcommand{\footrulewidth}{0.1pt}
\newcommand{\minipagewidth}{0.8\columnwidth}

\renewcommand{\thefootnote}{\fnsymbol{footnote}}


\renewcommand{\headrulewidth}{0.1pt}
\renewcommand{\footrulewidth}{0.1pt}

\renewcommand{\thefootnote}{\fnsymbol{footnote}}




\begin{document}


\title{\textsc{\hw\\
\coursenumber}}
\author{Instructor: \instructorname}
\date{Due: \duedate\\ 125 Points}

\maketitle

\setlength{\unitlength}{1in}
You are allowed to use a generative model-based AI tool for your assignment. However, you must submit an accompanying reflection report detailing how you used the AI tool, the specific query you made, and how it improved your understanding of the subject. You are also required to submit screenshots of your conversation with any large language model (LLM) or equivalent conversational AI, clearly showing the prompts and your login avatar. Some conversational AIs provide a way to share a conversation link, and such a link is desirable for authenticity. Failure to do so may result in actions taken in compliance with the plagiarism policy.

Additionally, you must include your thoughts on how you would approach the assignment if such a tool were not available. Failure to provide a reflection report for every assignment where an AI tool is used may result in a penalty, and subsequent actions will be taken in line with the plagiarism policy.

\subsection*{Submission instruction:}
You may either complete your entire homework on Python Notebook  and submit a Google Colab link or you may choose a combination of a pdf and Google Colab Notebook. However, you must provide \textbf{publicly accessible} Google Colab URL on Canvas.

For .pdf on Canvas, follow the format  \texttt{CPE486586-LastFirst-HW-XX}. For example, if your name is Sam Wells, your file name should be \texttt{CPE486586-WellsSam-HW-XX.pdf}.

Your submission must contain your name and UAH Charger ID or your UAH email address.  Please number your pages as well.

\textbf{Please write down your unique keyword corresponding to your package and your Github account URL. Do not provide a link to your Google Colab notebook. Upload your Google colab notebook to your repo in a new folder called \textbf{HWNotebooks}. Provide your GitHub commit hash. Your submitted notebook must be fully executed with all cells outputs. Please make sure there are no error in your submitted notebook.}

\vskip 1em
\hrule
\vskip 1em


% {\LARGE \color{stormblue} \textbf{Python Practice-I}}



{{\LARGE \textbf{Theory}}}

\section{Automatic Differentiation and Expression Graph (10 Points)}

Consider a function $f(w, x, y, z) = wx + y^2 z$. Draw the expression graph corresponding to the given function and demonstrate how would you evaluate $\cfrac{\partial f}{\partial w}$ for $w=1$, $ x= 2$, $ y = 3$, and $z = 4$.

\section{Unconstrained Optimization (5 Points)}

\begin{enumerate}
\item Minimize $f(x) = x^2 - 4x + 4$, determine the critical point, and repor the minimum value of the function. \hfill \textbf{(5 Points)}

\item Optimize $f(x) = x^3 - 3x^2$. Find all the critical points possible. Which one of them are local minima and local maxima? What is the local minimum and maxmium values of $f(x)$? \hfill \textbf{(10 Points)}


\item Consider $f(x, y) = x^2 - y^2$. Compute its gradient $\nabla f$ and associated critical point. Is this critical point a maximum, a minimum, or a saddle point? \hfill \textbf{(5 Points)}

\end{enumerate}


\section{Constrained Optimization using Lagrange Multiplier (10 Points)}
Using method of Lagrange mulitplier, minimize $f(x, y, z) = x^2 + y^2 + z^2$ subject to $x + 2y + z = 6$.

\begin{enumerate}
\item Set up the Lagrangian.
\item Derive the system of equations.
\item Solve for critical points.
\item Classify extrema.
\item Explain the role of the constraint.
\end{enumerate}

\section{Optimizing a Function using Stochastic Gradient Descent (5 Points)}
Use the Stochastic Gradient Descent technique to calculate the optimal value of 
$w_1$ for the cost function shown in Equation (1). Use a learning rate 
$\eta = 0.1$ , and an initial value 
$w_1^{(0)} = 2.0$.
Write down the updated value of 
$w_1$ for 5 iterations using the formula for gradient descent. \newline
This question is \textbf{analytical}, no code is required.
\begin{equation}
    J(w_1) = (w_1-5)^4 + (w_1-3)^2 +  4
\end{equation}

How does it compare with optimizing $J$ using partial differentiation method?

\section{Linear Regression and Regularization (15 points)}
In class we derived and discussed linear regression in detail. Find the result of minimizing the
loss of sum of the squared errors after adding in a penalty called an $L_2$ penalty on the weights.
More formally
\begin{multiequation}
\argmin_\wbf \bigg\{ ||\ybf - \xbf \wbf||^2  + \lambda||\wbf||^2  \bigg\}
\end{multiequation}

which means we are interested in finding $\wbf$ that minimizes the overall expression. How does this change the solution to the original linear regression solution? What is the impact of adding in this penalty?

$|| \cdot ||$ is L2 norm.

\textbf{Hint:} First take the gradient with respect to $\wbf$ and then equate to zero you will obtain an expression for $\wbf$ which is the estimated $\hat{\wbf}$.

Also, remember in vector form the L2 norm $||v||$ can be written as $\vbf^\top \vbf$.

\section{Assumptions in Linear Regression (5 Points)}

Write all the assumptions while considering simple linear regression model.

\bigskip

\bigskip

{{\LARGE \textbf{Practice}}}

Note that you don't need to use Google Colab for submission.

\section{Computing Derivative  using PyTorch Autograd (5 Points)}

Consider a following function
\begin{multiequation}
  \tau  = 2 r \cdot \sqrt{\cfrac{r}{r_s}} \cdot \sqrt{1 - \cfrac{r_s}{r}}
\end{multiequation}

which is the time experienced by an observer traveling along a specific trajectory in curved spacetime. $r$ is  the radial coordinate in a spherical coordinate system and $r_s$ is the Schwarzschild radius.

Write a program in PyTorch to compute the derivative $\cfrac{d\tau}{dr}$ considering the value of $r = 30000.1$  for $r_s = 0.00887$.
  
\section{Linear Regression in One Variable with Scikit-Learn (35 Points)}


Consider the dataset \url{https://github.com/rahulbhadani/Datasets/blob/master/RAV4/2020-04-15-05-12-00_2T3Y1RFV8KC014025_CAN_Messages_5_Hz.csv} that contains various driving features from a test driving.

In your Google Colab notebook, use Scikit-learn to create a linear regression model with steer angle as independent  variable and yaw rate as response variable. You need to split the dataset into training and test with 25\% of the rows set aside as a test data. Report the estimated coefficients of the fitted linear model and $R^2$ computed using the test data.



\section{Linear Regression in One Variable with PyTorch (35 Points)}

Consider the dataset \url{https://github.com/rahulbhadani/Datasets/blob/master/Hydropower.csv}. The main task as a part of this assignment is to create a linear model $y  = w_0 + w_1 x$ and determine the suitable parameter $w_0$ and $w_1$.  For this consider the independent variable $x$ as Benefit-Cost-Ratio (BCR) from the given dataset and the response variable $y$ as AnnualProduction. Using the linear model, we should be able to predict Annual Production given a new BCR.



You need to write LinearRegression class based on gradient descent implemented using PyTorch as a part of your Python package with the following specification:

\begin{enumerate}

\item Create a subpackage \textbf{model} in your UV Python package. \hfill \textbf{(2 Points)}

\item The subpackage should contain a module \texttt{regression.py}. \hfill \textbf{(2 Points)}

\item Add a \texttt{LinearRegression} class to the module with the following requirement:

\begin{enumerate}

\item \texttt{\_\_init\_\_} function, which is a constructor should  take following parameters as argument: (i) learning rate, and (ii) number of epochs. Further, \texttt{\_\_init\_\_} should initialize $w_0$ and $w_1$, SGD optimizer, and mean-squared error as a loss function, and list or vector to store intermediate $w_0$, $w_1$, and loss as the training progresses. \hfill \textbf{(5 Points)}

\item \texttt{forward} function should specify the linear model and compute the response $y$ and return it. \hfill \textbf{(2 Points)}

\item \texttt{fit} function should perform machine learning on training dataset provided to the function as an argument. fit function also should take test data as an additional argument and after training, it should compute $R^2$ on the test data and print it. \hfill \textbf{(5 Points)}

\item \texttt{predict} function should predict AnnualProduction given a new BCR data.  \textbf{(4 Points)}

\item \texttt{analysis\_plot} function should create several plots on the same figure using subplots such as original data, fitted regression line, nd $w_0$ , $w_1$ and loss as the training progressed. \hfill \textbf{(5 Points)}



\end{enumerate}


\item Add appropriate entry to \texttt{\_\_init\_\_.py} at subpackage level and at package level to facilitate proper import. \hfill \textbf{(2 Points)}

\item Build your package after incrementing version number, commit to GitHub, and publish to PYPI. \hfill \textbf{(3 Points)}

\item Show the complete usage of your LinearRegression model with \texttt{Hydropower.csv} dataset through a standalone notebook that first installs your package. \hfill \textbf{(5 Points)}
\end{enumerate}

\end{document}